% ---------------------------------------------------------------------
%  Chapitre 5 — Le protocole, ses limites, ses suites
%  Rôle : livrer L'OUTIL (le produit du stage), assumer ses limites,
%  et tracer les étapes suivantes (dont les chantiers retirés du
%  périmètre : logistique sur grappes, signaux d'information, lead-lag).
% ---------------------------------------------------------------------
\chapter{Du jugement à l'outil~: le protocole, ses limites, ses suites}
\label{chap:protocole}

\section{Le protocole récapitulé}
\label{sec:protocole}

% [À RÉDIGER — présenter l'outil comme le LIVRABLE du stage, 5-8 lignes]
% Idées :
%  - Ce que le stage livre n'est pas une stratégie gagnante mais une
%    procédure de jugement, réutilisable telle quelle sur tout signal
%    futur de l'agent (technique, informationnel ou relationnel).
%  - Chaque étage répare une condition d'application précise ; le
%    verdict n'est prononcé que si TOUS les étages concordent.
%  - Renvoyer au tableau récapitulatif.
Rappelons que le protocole ne fournit pas une stratégie gagnante, mais une
procédure de jugement, réutilisable telle quelle sur tout signal futur de
l'agent, qu'il soit technique, informationnel ou relationnel. Chaque étage
répare une condition d'application précise, et le verdict n'est prononcé que
si tous les étages concordent. Le tableau~\ref{tab:protocole} récapitule cette
procédure en six étages.

\begin{table}[H]
\centering
\caption{Le protocole de validation en six étages~: chaque étage répare
une menace identifiée, et le verdict exige de tous les franchir.}
\label{tab:protocole}
\small
\begin{tabular}{@{}p{0.30\linewidth}p{0.34\linewidth}p{0.26\linewidth}@{}}
\toprule
Étage & Menace traitée & Outil \\
\midrule
1. \edgevar{} apparié au hasard & dérive du marché & Monte Carlo apparié \\
2. Student unilatéral + TCL & variabilité d'échantillonnage & test de référence \\
3. Seuil de Bonferroni & multiplicité (11 tests) & contrôle du risque global \\
4. DEFF + bootstrap par grappes & dépendance intra-titre & sondages, rééchantillonnage \\
5. Les deux portes (titre, mois) & pseudo-réplication & agrégation équipondérée \\
6. Série mensuelle AR($p$) & autocorrélation résiduelle & variance de long terme \\
\bottomrule
\end{tabular}
\end{table}


\section{Limites assumées}
\label{sec:limites}

Aucune de ces limites n'invalide le verdict rendu~; chacune trace en revanche
la frontière de ce que le protocole peut affirmer. On les regroupe en limites
techniques, propres à la construction des données, et en limites statistiques,
propres aux méthodes elles-mêmes.

\subsection*{Limites techniques}

La première tient à la façon dont les stratégies ont été établies. Si l'on
étudiait des stratégies dont les règles ont été optimisées sur les données
mêmes qui servent ensuite à les juger, on tomberait dans le sur-apprentissage.
Le risque est ici atténué, car les stratégies retenues sont des règles
classiques à peu de paramètres, non ajustées sur l'échantillon~; il n'est pas
pour autant nul, et une validation hors échantillon, où l'on fige les règles
sur une période puis on juge sur la suivante, resterait nécessaire avant tout
usage opérationnel.

Par ailleurs, les frais et les conditions d'exécution n'ont été modélisés que
de façon simplifiée~; les intégrer plus finement rapprocherait le test des
conditions réelles.

Enfin, la base ne contient que les sociétés présentes dans l'indice sur la
période, les entreprises disparues en étant absentes. Ce biais du survivant
est atténué par la construction appariée de l'\edgevar{}, qui compare chaque
trade au hasard sur le même titre, mais il subsiste.

\subsection*{Limites statistiques}

Le test de Student initial concède beaucoup d'incertitude, puisqu'il surestime
la quantité d'information disponible en traitant les trades comme indépendants.
Il a néanmoins suffi à écarter huit stratégies. Il suppose aussi un effectif
assez grand pour invoquer le théorème central limite, condition d'autant plus
importante qu'en finance les données ne suivent guère une loi normale.

Le contrôle du risque global par la correction de Bonferroni est adapté à un
nombre restreint de stratégies déclarées à l'avance. Sur un univers de milliers
de variantes essayées, le contrôle de la fouille de données exigerait plutôt un
outil dédié comme le \emph{Reality Check} de White~\cite{white2000}.

Une limite plus subtile tient à l'empilement des corrections. Chaque méthode a
été conçue pour une seule source de dépendance~: le DEFF et le bootstrap pour
la dépendance intra-titre, chaque porte pour sa propre dimension. Prises
isolément, elles sont exactes~; utilisées ensemble sur des données où les
dépendances coexistent, elles perdent la garantie de leur seuil nominal, et
l'inférence de la correction mensuelle en particulier n'est plus exactement
calibrée. Le protocole reste cependant fixe et applique toujours les mêmes
étapes~; ce qui varie, c'est la lecture. Lorsqu'une dépendance se révèle
négligeable sur un canal, le seuil nominal y demeure valide et la conclusion
peut être chiffrée~; ce n'est que si les trois canaux se manifestaient
fortement en même temps, leurs corrections se cumulant, que le risque global
cesserait d'être chiffrable, et ce cas serait alors signalé. Comme la
dépendance simultanée aux trois canaux reste faible dans nos données, on
retient les verdicts obtenus.

Il faut enfin rappeler que la puissance élevée des grands échantillons rend le
moindre écart significatif~: significativité et taille d'effet doivent être
distinguées avec soin.

Face à chacune de ces limites, le protocole ne tranche pas au forceps~: il
préfère écarter une stratégie en signalant la réserve plutôt que de conclure
sur une base fragile.

\section{Perspectives}
\label{sec:perspectives}

% [À RÉDIGER — 3 paragraphes courts, un par perspective] Idées :
%  1. VALIDATION TEMPORELLE : figer les seuils sur une période
%     d'apprentissage, juger sur la suivante (walk-forward /
%     TimeSeriesSplit) — la validation croisée classique est interdite
%     ici (fuite du futur).
%  2. L'EDGE CONDITIONNEL : modéliser P(win | contexte) par une
%     régression logistique. Écartée de ce mémoire PAR RIGUEUR : sur
%     38 210 trades dépendants, les p-values des coefficients seraient
%     trop optimistes (même piège qu'au chap. 3). Les outils adaptés
%     (modèles à effets mixtes, équations d'estimation généralisées)
%     dépassent le cadre du M1 -> prochaine étape naturelle.
%  (bis) GARDE-FOU D'EFFECTIF : le protocole s'appuie sur le TCL, qui
%     exige un n suffisant. Ici tous les n dépassent 978, mais un outil
%     GÉNÉRIQUE devrait écarter en amont les stratégies trop peu
%     fournies (ni Shapiro ni le TCL ne permettent alors de conclure).
%     Ajout simple, à prévoir dans une version outillée du protocole.
%  3. LES AUTRES SIGNAUX DE L'AGENT : l'agent collecte aussi des
%     signaux d'information (contrats publics, transactions du
%     Congrès, réglementation) et des relations d'avance/retard entre
%     titres (causalité de Granger \cite{granger1969}). Des travaux
%     exploratoires ont été engagés sur ces deux fronts ; le protocole
%     du chap. 3 est précisément l'outil qui permettra de les juger —
%     hors du périmètre de ce mémoire.

La première perspective est la création de nos propres stratégies. Elle
suppose de séparer une zone d'apprentissage et une zone de test, en avançant
période après période plutôt que par une validation croisée classique, qui
mélangerait le futur et le passé. Le protocole devra être adapté en
conséquence.

Il faudra aussi lui ajouter un garde-fou d'effectif, écartant en amont les
stratégies trop peu fournies pour que le théorème central limite s'applique,
car ni le test de normalité ni le théorème central limite ne permettent alors
de conclure.

Une autre voie consiste à modéliser l'avantage en fonction du contexte, par
exemple la probabilité de gain selon le régime de marché ou le secteur, ce qui
prolongerait directement la piste des stratégies ciblées évoquée au
chapitre~\ref{chap:juger}. Cette voie a été écartée du présent travail par
rigueur~: sur des dizaines de milliers de trades dépendants, les p-values des
coefficients d'une régression seraient trop optimistes, exactement le piège du
chapitre~\ref{chap:juger}. Les outils adaptés, modèles à effets mixtes ou
équations d'estimation généralisées, dépassent le cadre de ce travail.

L'objectif est enfin de fournir à l'agent un maximum d'information, y compris
des signaux tournés vers l'avenir comme les projets de loi ou les signatures de
contrats, ainsi que les relations d'avance et de retard entre titres, au sens
de la causalité de Granger~\cite{granger1969}. Des travaux exploratoires ont
été engagés sur ces fronts. Ce mémoire en constitue la fondation~: il fournit
l'outil qui permettra, tout du long, de vérifier les stratégies à venir.